% WHAT A \valign PUTS BETWEEN THE COLUMNS ONE ENTRY SPANS.
%
% This is trip's line 337 cut down to its own file. That line says
%
%   A&\omit\valign to -5pt{#&#\cr A\char`}\span\cr{ }\span\cr}\cr
%
% -- a \valign set to a NEGATIVE height whose every row ends in \span, so
% every row covers both of the columns and no row ever starts in the second.
% What trip's second pass draws of it, twice over:
%
%   ref:    ....\glue(\tabskip) 0.0
%           ....\vbox(8.0+0.0)x0.0
%   hstex:  ....\glue(\tabskip) 4.49998 plus 7227.0
%           ....\vbox(3.50002+0.0)x0.0
%
% and 8.0 - 4.49998 = 3.50002 exactly, so the two are not disagreeing about
% any length: they are splitting one length between the glue and the box in
% two different ways. [A] below is the whole of that in six lines.
%
% WHAT THE REFERENCE DOES, measured over three \tabskip values -- 4.5pt,
% 10pt and nothing at all. The glue in front of the covered column is 0.0
% every time, and the column keeps the entry's whole extent, 6.83331, every
% time. HSTeX draws the \tabskip and gives the column what is left, which is
% 2.33331 at 4.5pt and -3.16669 at 10pt. Where it comes from is the
% measuring loop in finish_alignment: a span's excess goes to the LAST
% column it covers, and what it already covers is counted as the columns
% plus THE TABSKIPS BETWEEN THEM.
%
% THE RULE IS NOT `A \valign HAS NO GLUE INSIDE A SPAN'. That was tried --
% not counting the tabskip while measuring, and drawing zero in its place --
% and [B] is the case that refutes it: one row that spans and one that does
% not. There the reference draws the \tabskip in full, 10.0, in the spanning
% row as well, and the change made HSTeX draw 0.0. Backed out.
%
% [B] does not settle the measuring question, only the drawing one: its
% second row sizes both columns on its own, so the span has nothing left to
% make up and the loop never reaches the arithmetic that counts the tabskip.
%
% WHERE TO LOOK NEXT. What is different about [A] is that NO ROW EVER STARTS
% IN THE SECOND COLUMN. This engine already has a notion of how many columns
% a set of rows `reached', and drops the tabskip after the ones past it --
% see finish_alignment, `the ones past the last one used stand in the row the
% preamble is packed as, with no width and NO TABSKIP AFTER THEM'. A column
% that is only ever spanned INTO may be the same kind of thing, which would
% put the zero in the column's own record and explain the other place it
% shows: the packed preamble row of [A] draws `10.0 0.0 10.0' in the
% reference and `10.0 10.0 10.0' here, and one rule in the record would give
% both.
%
% An \halign is already exact for all of this; see
% what-an-entry-that-spans-is-set-to.tex, whose five cases agree box for box.
\catcode`\{=1 \catcode`\}=2 \catcode`\#=6 \catcode`\&=4
\tracingonline=1 \showboxbreadth=99 \showboxdepth=99
\font\tenrm=cmr10 \tenrm
\message{[A every row spans, which is trip line 337]}
\tabskip=4.5pt plus 7227pt
\setbox0=\hbox{\valign to -5pt{#&#\cr \hbox{A}\span\cr \hbox{B}\span\cr}}
\showbox0
\message{[A' the same with a plain tabskip, and with none]}
\tabskip=10pt
\setbox0=\hbox{\valign to -5pt{#&#\cr \hbox{A}\span\cr \hbox{B}\span\cr}}
\showbox0
\tabskip=0pt
\setbox0=\hbox{\valign to -5pt{#&#\cr \hbox{A}\span\cr \hbox{B}\span\cr}}
\showbox0
\message{[A'' three columns, every row spanning all of them]}
\tabskip=6pt
\setbox0=\hbox{\valign to -5pt{#&#&#\cr \hbox{A}\span\span\cr \hbox{B}\span\span\cr}}
\showbox0
\message{[B one row spans and one does not: the glue is drawn in full]}
\tabskip=10pt
\setbox0=\hbox{\valign to -5pt{#&#\cr \hbox{A}\span\cr \hbox{B}&\hbox{C}\cr}}
\showbox0
\message{[done]}
\end
